\scanpage{51}
\mspage{50}

que \unsure{pour} $i,j\in I$, $\unsure{i\Rightarrow j|i\ne j}$, $F'_i|F_i\cap\deleted{G}F_j=\deleted{G}F'_j|F_i\cap F_j$ Alors il existe une unique subdivision $F'$ de $F$ telle que $F'|F_i=F'_i\quad\forall i\in I$

\uline{Cor} Pour que deux subdivisions \inserted{$F',F''$ de $F$} soient égales, il \unsure{f} et \unsure{s} qu'elles induisent la $=$ subdivision sur chaque sous-figure \unsure{élémentaire} $F_{\unsure{X|x}}$ \unsure{.ie|i.e.} sur chaque \unsure{multistrate}.

\deleted{\unsure{Cor}} \uline{Proposition} \deleted{\unsure{1|a}}) Soient $F$ une figure, $G$ une sous-figure, $\deleted{G}F'$ un raffinement de $F$, et $G'=\{X\in F'\mid X\ll G\}$ c'est donc une sous figure de $F'$, \deleted{et} \inserted{Je dis qu'}elle raffine $G$ et c'est \inserted{\deleted{\unsure{d'ailleurs|même}}} la plus grande sous-figure de $F'$ qui raffine $G'$.

% slanted
\margin{\uline{raffinement} \uline{induit}. \uline{NB} Si $\deleted{G}F'\subset\deleted{G}F$, alors $G'$ n'est autre que l'intersection $F'\cap G'$.}

\diagram{\begin{tikzcd}[ampersand replacement=\&,column sep=small,row sep=small] \deleted{G}F' \arrow[r,phantom,"\ll" description] \arrow[d,phantom,"\cup" description] \& F \arrow[d,phantom,"\cup" description] \\ G' \arrow[r,phantom,"\ll" description] \& G \end{tikzcd}}

\deleted{\unsure{Ceci dit},} Il y a qque chose à prouver, car le seul fait que les $X\in G'$ raffinent $G$ n'implique \unsure{ps|pas} que $G'$ raffine $G$. Mais on a une \deleted{\unsure{raffinement}} \inserted{subdiv.} $\overline{F\unsure{''|'|{}}}$ de $F$ tel que \diagram{\begin{tikzcd}[ampersand replacement=\&,column sep=small,row sep=small] F' \arrow[r,phantom,"\subset" description] \& \overline{F''}\arrow[r,"\sim","s."' ]\arrow[d,phantom,"\cup" description] \& F\arrow[d,phantom,"\cup" description] \\ \& \overline{G''}\arrow[r,"\sim","s"' ] \& G \end{tikzcd}} \unsure{considérons|considérer} la subd. induite \deleted{sur} $G$\inserted{(\unsure{C,6|C6})}, $\overline{G''}$, formée des $X\in\overline{F}$ \unsure{telles|tels} que $X\ll \deleted{\unsure{F}}G$. \deleted{\unsure{Ceci}} Par définition, on a donc $G'=F'\cap G'$, qui est \deleted{bien} une sous-figure de $\overline{\deleted{\unsure{C}}G''}$ donc raffine $G$.
